% !PRINTED 415
\weberpage{415}
If one further sets
\weq{6}{a+\partial = y,}
then it follows:
\weq{7}{\begin{aligned}
a&=\frac{y+B\,x}{2},\\
b&=A\,x,\\
c&=-C\,x,\\
\partial&=\frac{y-B\,x}{2},
\end{aligned}}
from which, if
\[
a\,\partial - b\,c = n
\]
is set, there results:
\weq{8}{4\,n = y^2 - D\,x^2.}
The number $n$ is thus decomposed into the two complex integer factors
\[
n = \frac{y+x\sqrt{D}}{2}\cdot\frac{y-x\sqrt{D}}{2}
\].
For $\omega$ and $\mu$ there is found further:
\[
\omega = \frac{-B+\sqrt{D}}{2\,A} = \frac{2\,C}{-B-\sqrt{D}},
\]
\[
\mu = a+b\,\omega = \frac{y+\sqrt{D}\,x}{2}.
\]

The numbers $x$, $y$ introduced here are bound to the single condition
that $y+B\,x$, and consequently $y-B\,x$ as well, be even numbers, so
that $a$, $\partial$ become integers according to (7). If therefore $B$
is even, then $y$ must be even, while $x$ is arbitrary; if $B$ is odd,
then $x$ and $y$ are to be assumed both even or both odd. By (4) this
distinction can also be expressed thus:
\[
\begin{aligned}
\text{1) If }D&\equiv0\ (\mathrm{mod}\ 4), &\text{then }y&\equiv0\ (\mathrm{mod}\ 2),\\
\text{2) If }D&\equiv1\ (\mathrm{mod}\ 4), &\text{then }y&\equiv x\ (\mathrm{mod}\ 2).
\end{aligned}
\]

If, for given $A$, $B$, $C$, the numbers $a$, $b$, $c$, $\partial$ are to
be without a common divisor, then still other conditions come in:

From (5), (6) it follows that every common divisor of $a$, $b$, $c$,
$\partial$ is also a divisor of $x$ and $y$; conversely every common
divisor of $x$ and $y$ is also a divisor of $b$, $c$, $2\,a$,
$2\,\partial$. If therefore $a$, $b$, $c$, $\partial$ are to be without a
common divisor, then $x$ and $y$
%CONT>>
